\documentclass[11pt]{article}
\usepackage[margin=1in]{geometry}
\usepackage{amsmath,amssymb,amsthm}
\usepackage{booktabs}
\usepackage{algorithm}
\usepackage{algpseudocode}
\usepackage[colorlinks=true,linkcolor=blue,citecolor=blue,urlcolor=blue]{hyperref}
\usepackage{enumitem}

\newcommand{\y}{\mathbf{y}}
\newcommand{\x}{\mathbf{x}}
\newcommand{\g}{\mathbf{g}}
\newcommand{\rv}{\mathbf{r}}
\newcommand{\K}{\mathbf{K}}
\newcommand{\Mm}{\mathbf{M}}
\newcommand{\Am}{\mathcal{A}}
\newcommand{\X}{\mathbf{X}}
\newcommand{\Q}{\mathbf{Q}}
\newcommand{\U}{\mathbf{U}}
\newcommand{\V}{\mathbf{V}}
\newcommand{\G}{\mathbf{G}}
\newcommand{\Sm}{\mathbf{S}}
\newcommand{\R}{\mathbb{R}}
\newcommand{\tr}{\operatorname{tr}}
\newcommand{\E}{\mathbb{E}}
\newcommand{\given}[1]{\,\middle|\,#1}
\title{mixmogam 2.0: Technical Details of the Mixed-Model Association Engine}
\author{mixmogam contributors}
\date{October 2026}

\begin{document}
\maketitle

\section{Overview and notation}

mixmogam fits the standard linear mixed model for genome-wide
association mapping,
\begin{equation}
  \y = \X\beta + u + e, \qquad
  u \sim \mathcal{N}(0, \sigma_g^2 \K), \quad
  e \sim \mathcal{N}(0, \sigma_e^2 I_n),
  \label{eq:model}
\end{equation}
where $n$ is the number of samples, $\X \in \R^{n\times q}$ the
covariate design (an intercept is included by default), and
$\K \in \R^{n\times n}$ a kinship (relationship) matrix. We write
$\delta = \sigma_e^2/\sigma_g^2$ and
$\V = \sigma_g^2(\K + \delta I_n)$. The covariance of \eqref{eq:model}
is parameterized by the single ratio $\delta$; all variance-component
estimation reduces to one-dimensional optimization of the REML or ML
profile likelihood over $\delta$.

The package implements three regimes:
\begin{enumerate}[nosep]
  \item \textbf{Exact} ($n \lesssim 8{,}000$): one dense
        eigendecomposition of $\K$ and of the restricted matrix
        $\Sm(\K+I)\Sm$ (Section~\ref{sec:fit}); the association scan
        evaluates the exact EMMAX statistic in batched dense GEMMs
        (Section~\ref{sec:scan}).
  \item \textbf{Truncated spectrum} (large $n$): a randomized top-$k$
        eigenbasis of $\K$ with an exact identity-tail correction
        (Section~\ref{sec:topk}).
  \item \textbf{SLQ fitting} (large $n$): variance components by
        deflated stochastic Lanczos quadrature, requiring only
        matrix--vector products with $\K$ and no $O(n^3)$
        factorization (Section~\ref{sec:slq}).
\end{enumerate}
Genotypes are stored sample-major as $n \times m$ matrices of
\texttt{int8} dosages in $\{0,1,2\}$ with $-1$ encoding missing calls;
SNP blocks are converted to \texttt{float32} with fused mean
imputation immediately before use, so the float working set is bounded
by the block size regardless of $m$.

\section{Variance-component estimation: exact EMMA}
\label{sec:fit}

\subsection{Restricted eigenspace}

Let $\Sm = I_n - \X(\X^\top\X)^{-1}\X^\top$ project onto the orthogonal
complement of the covariate space. Following
Kang et al.\ (2008), the REML likelihood depends on $\y$ and $\K$ only
through the spectral decomposition of the \emph{restricted} matrix
\begin{equation}
  \Mm = \Sm(\K + I_n)\Sm .
\end{equation}
Writing its eigenpairs $(\mu_i, u_i)$ in ascending order, the first
$q$ eigenvalues are (numerically) zero; for the remaining $p = n - q$
directions we set
\begin{equation}
  \lambda^R_i = \mu_{q+i} - 1, \qquad
  \eta_i = u_{q+i}^\top \y .
\end{equation}
The factorization is computed once in \texttt{float64} and cached.

\subsection{Profile likelihood and grid}

The REML profile log-likelihood as a function of $\delta$ is
\begin{equation}
  \ell_R(\delta) = \tfrac12\Big[p\big(\log\tfrac{p}{2\pi} - 1 - \log s_1(\delta)\big) - s_2(\delta)\Big],
  \qquad
  s_1(\delta) = \sum_{i=1}^{p} \frac{\eta_i^2}{\lambda^R_i + \delta},
  \quad
  s_2(\delta) = \sum_{i=1}^{p} \log(\lambda^R_i + \delta),
  \label{eq:reml}
\end{equation}
and the ML variant replaces $p$ by $n$ and $s_2$ by the full-spectrum
determinant $\sum_{j=1}^{n}\log(\lambda^K_j + \delta)$ evaluated on the
eigenvalues of $\K$, using the R~EMMA formulation (the quadratic form
$s_1$ retains the restricted basis). The grid
$\delta_j = \exp(\mathrm{linspace}(l_{\min}, l_{\max}, G{+}1))$ with
$G = 100$ and $[l_{\min},l_{\max}] = [-10, 10]$ is evaluated
\emph{vectorized over all grid points simultaneously} by broadcasting
$\lambda^R$ against $\delta$, so the whole grid costs $O(pG)$ fused
operations. The maximizing bracket is located by a sign change of
$\mathrm{d}\ell/\mathrm{d}\delta$ and polished with Brent's method
(v1 used Newton from the bracket midpoint; the root is identical and
Brent cannot diverge), with v1's boundary-acceptance and
grid-argmax fallback rules retained. Finally
\begin{equation}
  \hat\sigma_g^2 = \frac{s_1(\hat\delta)}{p}, \qquad
  \hat\sigma_e^2 = \hat\delta\,\hat\sigma_g^2, \qquad
  h^2_{\text{pseudo}} = \frac{1}{1+\hat\delta}.
\end{equation}

\begin{algorithm}[t]
\caption{EMMA exact fit (REML shown; ML analogous)}
\label{alg:emma}
\begin{algorithmic}[1]
\Require $\y$, $\X$, $\K$, grid $(G, l_{\min}, l_{\max})$, tolerance $\tau$
\State $\Q \gets$ thin QR of $\X$;\quad $\Sm \gets I - \Q\Q^\top$
\State eigendecompose $\Mm = \Sm(\K + I)\Sm$; keep $\{\lambda^R_i, \eta_i\}_{i\le p}$
\State vectorize $\ell_R(\delta_j)$, $\ell'_R(\delta_j)$ over the grid $j=0..G$
\State $i^\ast \gets \arg\max_j \ell_R(\delta_j)$
\State find $i$ with $\ell'_R(\delta_{i}) > 0 > \ell'_R(\delta_{i+1})$ maximizing mean endpoint likelihood
\State $\hat\delta \gets \mathrm{Brent}(\ell'_R,\ \delta_i,\ \delta_{i+1},\ \tau)$; accept per boundary rules, else fall back to $\delta_{i^\ast}$
\State \Return $\hat\delta$, $\hat\sigma_g^2 = s_1(\hat\delta)/p$, $h^2_{\text{pseudo}}$, $\beta$ via \eqref{eq:gls}
\end{algorithmic}
\end{algorithm}

\subsection{GLS coefficients and gBLUP}

With $\H \equiv \V^{-1/2}$ (applied implicitly, Section~\ref{sec:topk}),
\begin{equation}
  \hat\beta = \big(\X^\top\V^{-1}\X\big)^{-1}\X^\top\V^{-1}\y
  \quad\text{computed as}\quad
  \big((\H\X)^\top(\H\X)\big)^{-1}(\H\X)^\top(\H\y),
  \label{eq:gls}
\end{equation}
and the gBLUP of breeding values is
$\hat u = \hat\sigma_g^2\,\K\,\V^{-1}(\y - \X\hat\beta)$, with
$\V^{-1}$ applied through the eigenbasis.

\section{The association scan: EMMAX as batched GEMM}
\label{sec:scan}

\subsection{Single-SNP statistic}

For a candidate SNP $g \in \R^n$ (missing calls mean-imputed), v1
regressed the transformed, covariate-residualized phenotype on the
transformed, covariate-residualized SNP by an explicit least-squares
solve \emph{per SNP}. The identical statistic has a closed form.
Define, at the fitted $\hat\delta$,
\begin{equation}
  \Q \gets \operatorname{qr}\big(\H\X\big), \qquad
  r \gets (I - \Q\Q^\top)\,\H\y, \qquad
  \tilde g \gets (I - \Q\Q^\top)\,\H g, \qquad
  \mathrm{RSS}_0 \gets \|r\|^2 .
\end{equation}
The one-degree-of-freedom F statistic is the rank-one update
\begin{equation}
  \mathrm{RSS}(g) = \mathrm{RSS}_0 - \frac{(\tilde g^\top r)^2}{\tilde g^\top\tilde g},
  \qquad
  F(g) = \Big(\frac{\mathrm{RSS}_0}{\mathrm{RSS}(g)} - 1\Big)(n - q - 1),
  \qquad
  p(g) = F_{1,\,n-q-1}\big(F(g)\big),
  \label{eq:scanstat}
\end{equation}
which equals v1's per-SNP regression exactly (Frobenius--Waugh--Lovell
partialling out). Effect sizes and standard errors come in closed form,
\begin{equation}
  \hat\beta_g = \frac{\tilde g^\top r}{\tilde g^\top \tilde g},
  \qquad
  \operatorname{se}(\hat\beta_g) =
  \sqrt{\frac{\mathrm{RSS}(g)}{(n - q - 1)\,\tilde g^\top\tilde g}} .
\end{equation}

\subsection{Batching}

All per-SNP quantities are rows of matrix expressions. For a block
$\G \in \R^{k\times n}$ of $k$ SNPs:
\begin{equation}
  \widetilde\G^\top = (I - \Q\Q^\top)\,\H\,\G^\top, \qquad
  \mathbf{num} = \widetilde\G\,^\top r \in \R^k, \qquad
  \mathbf{den} = \operatorname{diag}(\widetilde\G^\top\widetilde\G) \in \R^k,
\end{equation}
and \eqref{eq:scanstat} is applied elementwise. The transform
$\H\,\G^\top$ is two GEMMs against the eigenbasis (one GEMM per side);
the residualization is one small GEMM of cost $O(nqk)$. The scan is
therefore $O(n^2 m)$ dense linear algebra with BLAS-3 kernels, no
Python-level per-SNP loop, and no per-SNP factorization; arithmetic
defaults to \texttt{float32} (fitting stays \texttt{float64}).

\begin{algorithm}[t]
\caption{Batched EMMAX scan}
\label{alg:scan}
\begin{algorithmic}[1]
\Require fitted $\hat\delta$; covariates $\X$; phenotype $\y$; genotype store returning blocks $\G \in \R^{k\times n}$; block size $k$, dtype
\State $\Q \gets \operatorname{qr}(\H\X)$; $r \gets (I-\Q\Q^\top)\H\y$; $\mathrm{RSS}_0 \gets \|r\|^2$; $\nu \gets n - q - 1$
\ForAll{blocks $\G$ from the store}
  \State $\widetilde\G^\top \gets (I - \Q\Q^\top)\,\H\,\G^\top$ \Comment{2 GEMMs vs.\ eigenbasis + 1 small GEMM}
  \State $\mathbf{num} \gets \widetilde\G^\top r$;\quad $\mathbf{den} \gets \operatorname{rowsums}(\widetilde\G \circ \widetilde\G)$
  \State $t^2 \gets \mathbf{num}^2 / \max(\mathbf{den}, \epsilon)$;\quad
         $\mathbf{RSS} \gets \max(\mathrm{RSS}_0 - t^2, 0)$
  \State $F \gets t^2 / \mathbf{RSS} \cdot \nu$;\quad $p \gets F_{1,\nu}(F)$ \Comment{batched survival function}
\EndFor
\State \Return $\{p, F, \mathrm{RSS}, \hat\beta_g, \operatorname{se}(\hat\beta_g)\}$
\end{algorithmic}
\end{algorithm}

\subsection{Multi-df extensions}

2-df genotypic tests (v1's ANOVA path) group SNPs by level count $L$
and build $L{-}1$ indicator contrasts per SNP; the small
$(L{-}1)\times(L{-}1)$ normal equations are solved \emph{batched} over
the block (stacked \texttt{solve}). GxE scans test the interaction
column $g \circ E$ (1~df), or jointly $[g,\; g\circ E]$ (2~df) with the
same batched two-column update. The stepwise MLMM rescans with the
current cofactors appended to $\X$, selecting by
$\mathrm{BIC} = \ell - \tfrac{k}{2}\log n$ with backward elimination.

\section{Large \texorpdfstring{$n$}{n}: the kinship as a streaming operator}
\label{sec:operator}

Following BOLT-LMM and LDAK-KVIK, above the dense regime the kinship is
never materialized. \texttt{GenotypeKinship} wraps the genotype store and
applies $K$ only through streaming products,
\begin{equation}
  K X \;=\; \frac{1}{m}\sum_{\text{blocks}} Z_b \bigl(Z_b^\top X\bigr),
\end{equation}
with $Z_b$ the Yang-2010 called-only standardized block; the
\texttt{scale\_k} normalization constants (mean off-diagonal, mean
diagonal) are obtained from two cheap streaming passes and folded into
the same products, so the operator equals the dense GRM to machine
precision (verified to $10^{-14}$ in the tests). Every large-$n$
computation is operator-native: the randomized eigensolver, the SLQ
solver of Section~\ref{sec:slq}, the scan transform (through the
truncated basis below) and gBLUP all consume mat-vec products only.
Cost per product is $O(nm)$ tall-skinny GEMMs with $O(\text{block})$
working memory, replacing $O(n^2)$ storage and an $O(n^2 m)$
materialization; the equivalence to the dense pipeline (variance
components, likelihood, gBLUP, scan under gentle truncation) is tested
directly.

\paragraph{Adaptive deflation.} The SLQ solver deflates only
\emph{well-separated} extreme eigenvalues (Ritz values at least
$4\times$ the bulk median): the analytic accounting for deflated
directions is exact only for an invariant subspace, and subspace
iteration cannot converge the boundary of a near-degenerate bulk.
Spectra without outliers deflate nothing -- the tight bulk is precisely
the regime Gauss quadrature resolves well.

\section{Large \texorpdfstring{$n$}{n}: truncated randomized spectrum}
\label{sec:topk}

Above $n \approx 8{,}000$ (or whenever the kinship is a streaming
operator) the dense eigendecomposition is replaced by a randomized
top-$k$ approximation ($k = 2048$ by default). A Gaussian
sketch $\Omega \in \R^{n\times(k+p)}$ ($p = 12$ oversampling) is
powered by $\ell = 5$ subspace iterations,
$\Q_0 \gets \operatorname{orth}(\K\Omega)$, \
$\Q_t \gets \operatorname{orth}(\K\Q_{t-1})$, and the Ritz extraction
$B = \Q^\top\K\Q$, $B = W\Lambda W^\top$ yields
$\hat\Lambda_k, \hat U_k = \Q W_k$ (Halko et al.\ 2011). The
$V^{-1/2}$ map applied to any matrix $A$ becomes
\begin{equation}
  \boxed{\;
  \H A \;=\; \delta^{-1/2} A \;+\;
  \hat U_k\Big[\operatorname{diag}\big((\hat\lambda_i + \delta)^{-1/2}\big) -
  \delta^{-1/2} I_k\Big]\hat U_k^\top A \;}
  \label{eq:tail}
\end{equation}
which is \emph{exact} whenever the discarded eigenvalues vanish, costs
two $n \times k$ GEMMs, and never forms an $n \times n$ matrix. The
unexplained spectral mass
$\tau = \tr(\K) - \sum_{i\le k}\hat\lambda_i$ is reported on the fit so
the approximation can be audited; the same identity serves the scan,
gBLUP and $\V^{-1}$ applications.

\section{Large-\texorpdfstring{$n$}{n} fitting: deflated stochastic Lanczos quadrature}
\label{sec:slq}

Fitting $\delta$ without the full spectrum requires the spectral sums
$s_1, s_2, s_3, s_4$ of Section~\ref{sec:fit}. The SLQ solver
estimates them from Gauss quadrature rules built by Lanczos iteration
(Ubaru et al.\ 2017): with $T_t$ the tridiagonal Lanczos matrix after
$t$ steps from start vector $v$,
\begin{equation}
  v^\top f(\Am)\,v \;\approx\; \|v\|^2 \sum_{i=1}^{t} w_i f(\theta_i),
  \qquad (\theta, w) = \text{eigenpairs of } T_t,\ w_i = (S_{1i})^2,
\end{equation}
exact for polynomials of degree $\le 2t{-}1$ and spectrally accurate
for analytic $f$ such as $f_\delta(t) = (t + \delta)^{-1}$ and
$\log(t+\delta)$. Trace terms use unbiased Rademacher probe averages
($\E[z^\top f(\Am) z] = \tr f(\Am)$, probes pre-projected into the
covariate complement for REML); the quadratic forms in $\y$ use one
deterministic rule from the projected phenotype. All Lanczos runs use
full reorthogonalization ($t \le 96$ default, 12 probes).

\paragraph{Deflation.} Gauss quadrature converges slowly when the
spectrum mixes extreme outliers (population structure) with a tight
bulk. Extreme eigenpairs of the operator are therefore estimated by a
randomized top-$d$ pass ($d = 128$) and handled analytically; the
quadrature resolves only the deflated bulk. Because the deflated
operator carries the removed directions as (near-)zero eigenvalues,
the probe estimates must be corrected by their contribution
$f_\delta(0)$:
\begin{equation}
  \sum_{i>p-d} f(\lambda_i) \;=\;
  \Big[\text{probe mean of } f_\delta\text{-rules}\Big] - d\,f_\delta(0),
  \qquad f_\delta(0) \in \{1/\delta,\ \log\delta,\ \delta^{-2}\}.
\end{equation}
Quadrature nodes are floored at zero to guard against roundoff-negative
deflated eigenvalues. On structure-rich kinships the resulting
$\hat\delta$ agrees with the exact EMMA fit to a few percent (the REML
profile is very flat near its optimum) and profile likelihoods agree to
about one nat; the total cost is $O(n^2)$ per matvec times
$O(\text{probes} \times \text{steps})$ matvecs, with no $O(n^3)$
factorization.

\begin{algorithm}[t]
\caption{Deflated SLQ variance-component fit}
\label{alg:slq}
\begin{algorithmic}[1]
\Require operator $\mathrm{mv}(x) = \Sm\K\Sm x$; $\y$, $\X$; probes $P$, steps $t$, deflation $d$
\State $\hat\lambda, \hat U \gets$ randomized top-$d$ of $\mathrm{mv}$
\State $\mathrm{mv}_d(x) \gets \mathrm{mv}(x) - \hat U(\hat\lambda \circ (\hat U^\top x))$
\State $y_\perp \gets \Sm\y$; decompose $y_\perp = \hat U c + y_{\mathrm{rest}}$
\State rules: one Lanczos$(\mathrm{mv}_d, y_{\mathrm{rest}}, t)$; $P$ Lanczos$(\mathrm{mv}_d, \Sm z, t)$
\For{$\delta$ on grid and within Brent bracket}
  \State $s_1(\delta) \gets \sum_i c_i^2/(\hat\lambda_i{+}\delta) + \mathcal{Q}_y[f_\delta]$
  \State $s_2(\delta) \gets \sum_i \log(\hat\lambda_i{+}\delta) + \overline{\mathcal{Q}}_P[\log] - d\log\delta$ \Comment{deflation correction}
  \State $s_4(\delta) \gets \sum_i (\hat\lambda_i{+}\delta)^{-1} + \overline{\mathcal{Q}}_P[f_\delta] - d/\delta$
  \State assemble $\ell_R(\delta)$, $\ell'_R(\delta)$ as in \eqref{eq:reml}
\EndFor
\State \Return $\hat\delta$ by bracketing + Brent; $\hat\sigma_g^2 = s_1(\hat\delta)/p$
\end{algorithmic}
\end{algorithm}

\section{Kinships}

\paragraph{Additive GRM.} Per variant $j$, genotypes are standardized
over \emph{called} samples only,
$z_{ij} = (g_{ij} - \bar g_j)/s_j$ for called $g_{ij}$ and $z_{ij}=0$
otherwise (Yang et al.\ 2010), and
$\K = \sum_j w_j\, z_j z_j^\top / \sum_j w_j$ is accumulated over SNP
blocks in \texttt{float64}. Optional weights $w_j$ implement LDAK-style
weighting; an optional SNP subset fits the null model from a marker
subsample (LDAK-KVIK style) at reduced cost.

\paragraph{IBS.} Identity-by-state similarity is accumulated as one-hot
GEMMs per block: with $H_a$ the indicator of genotype $a \in \{0,1,2\}$
over called entries, the shared-allele numerator is
$\sum_a H_a H_a^\top$ and the denominator is $C C^\top$ with $C$ the
called-indicator, so $\K_{\mathrm{IBS}} = S/C$ per pair.

\paragraph{Exact LOCO.} Because the GRM is a sum of globally
standardized outer products, the blocks accumulated for each
chromosome $\K_c = m^{-1}\sum_{j \in c} z_j z_j^\top$ satisfy $\sum_c \K_c = \K$
\emph{exactly} under a single global standardization, and the
leave-one-chromosome-out matrix is the algebraic difference
\begin{equation}
  \K_{\mathrm{loco}}(c) = \frac{m\,\K - m_c\,\K_c}{m - m_c},
\end{equation}
rescaled to mean diagonal one. No per-chromosome restandardization is
performed, hence no approximation. Windowed local/global pairs along
the genome are the same construction restricted to a sliding interval.

\section{Batched permutation thresholds}

Permutation follows v1 semantics: the \emph{raw} phenotype is permuted
and variance components stay fixed at $\hat\delta$. Because
permutation does not commute with the whitening transform
($P\H \neq \H P$), permuted phenotypes are transformed after
permutation, $\mathbf{Y}_P = \H\,\y[\Pi]$, and residualized against the fixed
$\Q$. All $B$ permutations then flow through each SNP-block GEMM at
once:
\begin{equation}
  T = \widetilde\G^\top R_P \in \R^{k\times B}, \qquad
  F = \frac{T \circ T}{\mathbf{den}\,\mathrm{RSS}_{0,P}}\,\nu,
\end{equation}
so a permutation costs a matrix multiply rather than a full scan, and
the genome-wide 5\% threshold is the 5th percentile of the per-
permutation minimum $p$.

\section{Two-kinship mixtures}

For $\V = \sigma_{g_1}^2\K_1 + \sigma_{g_2}^2\K_2 + \sigma_e^2 I$, the
ratio $\rho = \sigma^2_{g_1}/\sigma^2_{g_2}$ is profiled on a grid of
$\log\rho$ over $[-3,3]$ (cascading refinements around the best
bracket), fitting each candidate mixture
$a\K_1 + (1{-}a)\K_2$ by the exact EMMA path of Section~\ref{sec:fit};
variance shares are reported per matrix.

\section{Missing data}

Per-SNP sample subsetting (v1's behavior) breaks block batching. The
default policy is a per-variant missingness filter followed by
mean imputation over called genotypes (EMMAX-consistent), fused into
the int8-to-float block conversion (optionally a Numba kernel);
\texttt{impute='none'} propagates NaN for analyses that prefer explicit
handling.

\section{Validation and measured performance}

Correctness is certified against a float64 port of the v1
mathematics (\texttt{tests/\_reference.py}): F statistics agree to
rtol $10^{-6}$, null $p$-values pass a Kolmogorov--Smirnov uniformity
test, simulated causal loci are recovered with high power, the SLQ
solver reproduces the exact fit (Section~\ref{sec:slq}), and the
bundled \emph{A.\ thaliana} Atwell data runs end-to-end. Measured
timings (development laptop, Apple silicon, 4 BLAS threads, AC power;
archive \texttt{benchmarks/results/}):

\begin{center}
\small
\setlength{\tabcolsep}{4.5pt}
\begin{tabular}{@{}llr@{}}
\toprule
Workload & Metric & Result \\
\midrule
scan vs.\ v1-style per-SNP least squares ($n{=}2000$, $m{=}10$k) & speed-up & $123\times$ (0.27\,s vs 33.7\,s) \\
scan, $n{=}500$, $m{=}100$k, float32 & throughput & $2.5\times10^{5}$ SNPs/s \\
scan, $n{=}2000$, $m{=}100$k, float32 & throughput & $4.8\times10^{4}$ SNPs/s ($3.3\times$ float64) \\
500 permutations $\times$ 5k SNPs, $n{=}1500$ (batched) & throughput & $1.7\times10^{6}$ perm-SNPs/s \\
\bottomrule
\end{tabular}
\end{center}

\section{References}

\begin{itemize}[nosep,leftmargin=1.4em]
\item Kang, H.M.\ et al.\ (2008) Efficient control of population
  structure in model organism association mapping. \emph{Genetics} 178:1709--1723. (EMMA)
\item Kang, H.M.\ et al.\ (2010) Variance component model to account
  for sample structure in genome-wide association studies.
  \emph{Nat Genet} 42:348--354. (EMMAX)
\item Yang, J.\ et al.\ (2010) Common SNPs explain a large proportion
  of the heritability for human height. \emph{Nat Genet} 42:565--569. (GRM)
\item Loh, P.-R.\ et al.\ (2015) Efficient Bayesian mixed-model
  analysis in genome-wide association studies. \emph{Nat Genet}
  47:284--290. (BOLT-LMM; truncated-spectrum scanning)
\item Mbatchou, J.\ et al.\ (2021) Computationally efficient
  whole-genome regression for quantitative and binary traits.
  \emph{Nat Genet} 53:1097--1103. (REGENIE; batched permutation-style projection scanning)
\item Tian, J.\ et al.\ (2024) LDAK-KVIK: fast and accurate
  variance-component estimation and association testing.
  (randomized eigensolvers, SNP-subsampled null fits)
\item Halko, N., Martinsson, P.-G., Tropp, J.A.\ (2011) Finding
  structure with randomness. \emph{SIAM Review} 53:217--288.
\item Ubaru, S., Chen, J., Saad, Y.\ (2017) Fast estimation of
  $\tr f(A)$ via stochastic Lanczos quadrature. \emph{SIAM J.\ Matrix Anal.\ Appl.}\ 38:1075--1099.
\end{itemize}

\end{document}
